\chapter{热力学基本概念与基本定律}

\begin{problem}
热力系的几种类型之间有何联系与区别？
\end{problem}

\begin{solution}
它们都是将所要研究的物质总和被人为划定边界后作为研究对象（即热力系），其核心区别在于系统边界对物质与能量（热量和功）交换的限制条件不同：

\begin{enumerate}
  \item \textbf{开口系}：与外界有物质交换和能量交换的系统。
  \item \textbf{封闭系（闭口系）}：与外界无物质交换的系统。
  \item \textbf{绝热系}：与外界无热量交换的系统。
  \item \textbf{孤立系}：与外界既无物质交换也无能量交换的系统。
\end{enumerate}
\end{solution}

\begin{problem}
工质的基本状态参数和导出状态参数分别是指哪些参数？
\end{problem}

\begin{solution}
\begin{enumerate}
  \item \textbf{基本状态参数}：指可以直接测量得到的状态参数，如温度（\(T\)）、压强（\(p\)）等。
  \item \textbf{导出状态参数}：指无法直接测量，只能用基本参数依某种关系推导而得的参数，如比热力学能（\(u\)）、比焓（\(h\)）、比熵（\(s\)）等。
\end{enumerate}
\end{solution}

\begin{problem}
试说明热力学第一定律和第二定律分别解决了热功转换及热量传递过程中的什么问题。
\end{problem}

\begin{solution}
\begin{enumerate}
  \item \textbf{热力学第一定律}：揭示了能量既可转换又要守恒，主要解决热变功过程的\textbf{数量计算}问题（即 energy 之间数量转换的关系）。
  \item \textbf{热力学第二定律}：揭示了能量转换具有方向性、条件和限度，主要解决热变功过程的\textbf{方向}问题（即自发过程不可逆，解释了能量不仅有数量守恒，更具有品位和转化方向）。
\end{enumerate}
\end{solution}

\begin{problem}
闭口系与开口系的第一定律表达式有何区别？
\end{problem}

\begin{solution}
这两种表达式的区别源于工质在流动过程中与外界交换能量的方式不同：

\begin{enumerate}
  \item \textbf{闭口系第一定律解析式}：\(q=w+\Delta u\)。表明加给工质的热量，一部分用来改变工质的内能，另一部分则用来使工质膨胀而对外作功。
  \item \textbf{开口系（稳定流动）第一定律解析式}：\(q=\Delta h+w_{t}\)。由于系统存在连续的物质进出，必须考虑流动带来的推挤功。因此引入了综合状态参数焓（\(h=u+p\upsilon\)），其物理图景变为：外界加入的热量，等于系统比焓的变化量与对外输出的技术功之和。
\end{enumerate}
\end{solution}

\begin{problem}
卡诺循环的意义是什么？
\end{problem}

\begin{solution}
\begin{enumerate}
  \item 从理论上确定了通过热机循环实现热能转变为机械能的条件，指出了提高热机热效率的方向，是研究热机性能不可缺少的准绳。
  \item 证明了卡诺循环热机效率仅取决于恒温热源的高温（\(T_{1}\)）和低温（\(T_{2}\)），而与工质的具体性质无关。
  \item 对热力学第二定律的建立具有重大意义。
\end{enumerate}
\end{solution}
